\section{Methodology}
\label{sec:methodology}

Building on the insight that policy change budget (KL divergence) should be
spent on structurally meaningful code transformations, we define structural
efficiency as a ratio and describe the four-component framework that operationalizes it.

\subsection{Problem Setup}

Let $\pi_\theta$ denote a code generation model fine-tuned from a base model
$\pi_\text{ref}$ using a post-training algorithm $A$ (e.g., GRPO or DPO).
Given a prompt $x$ and a generated completion $y$, define:

\begin{itemize}
  \item \textbf{KL divergence:} $D_\text{KL}(\pi_\theta \| \pi_\text{ref}) = \mathbb{E}_{x,y \sim \pi_\theta}[\log \pi_\theta(y|x) - \log \pi_\text{ref}(y|x)]$, measuring total policy divergence from the base model at checkpoint step $t$.
  \item \textbf{Semantic AST edit distance:} $d_\text{sem}(y, y')$, measuring minimum-cost tree edit operations restricted to control-flow and data-flow AST nodes between two code completions.
\end{itemize}

\textbf{Definition (Structural Efficiency):}

\begin{equation}
\text{SE}(\pi_\theta, \pi_\text{ref}) = \frac{\mathbb{E}[d_\text{sem}(y_\theta, y_\text{ref})]}{D_\text{KL}(\pi_\theta \| \pi_\text{ref})}
\label{eq:structural_efficiency}
\end{equation}

where $y_\theta \sim \pi_\theta(\cdot|x)$ and $y_\text{ref} \sim \pi_\text{ref}(\cdot|x)$
are completions from the fine-tuned and reference models respectively.

A high structural efficiency indicates that the method concentrates its policy
change budget on semantically meaningful code transformations. A low structural
efficiency indicates that the KL budget is consumed by surface-level changes
(variable naming, whitespace, comment style) that do not affect program behavior.

\subsection{FA-AST Node Classification}

\textbf{Rationale:} To restrict edit distance to semantically relevant nodes, we use
the FA-AST taxonomy \cite{Wang2020Detecting}, which classifies Python AST nodes into:

\begin{itemize}
  \item \textbf{CONTROL\_FLOW nodes:} \texttt{\{If, For, While, Try, With\}} --- nodes governing execution branching and looping
  \item \textbf{DATA\_FLOW nodes:} \texttt{\{Assign, Call, Return, FunctionDef\}} --- nodes governing data transformation and function structure
  \item \textbf{SURFACE nodes:} all other node types (constants, names, operators, formatting)
\end{itemize}

The Semantic Edit Proportion (SEP) measures what fraction of total edits targets
semantic (CF+DF) nodes:

\begin{equation}
\text{SEP} = \frac{d_\text{CF}(y_\theta, y_\text{ref}) + d_\text{DF}(y_\theta, y_\text{ref})}{d_\text{total}(y_\theta, y_\text{ref})}
\label{eq:sep}
\end{equation}

where $d_\text{CF}$, $d_\text{DF}$, $d_\text{total}$ are edit distances restricted
to control-flow nodes, data-flow nodes, and all nodes respectively.

Figure~\ref{fig:ast_node_heatmap} illustrates the FA-AST node classification
across checkpoint pairs, showing node-type frequencies for GRPO and DPO.

\subsection{ZSS Tree Edit Distance}

\textbf{Rationale:} AST-level edit distance captures structural program changes that
token-level metrics miss. Two code completions that differ only in variable names
are structurally identical at the AST level; two completions with different loop
structures are structurally distinct regardless of surface similarity.

We use the ZSS algorithm \cite{Zhang1989Simple}, which computes minimum-cost
tree edit distance in $O(n^2 m^2)$ time for trees of sizes $n$ and $m$. Edit
operations are node insertion, deletion, and relabeling, each with unit cost.
We restrict the edit distance computation to the semantic node types
defined in Section~\ref{sec:methodology}: an insertion or deletion of a control-flow or data-flow
node counts; an insertion of a \texttt{Name} node (surface) does not.

\textbf{Why not token-level metrics:} BLEU/CodeBLEU operate on token sequences and
are sensitive to variable renaming and formatting. Line-level diffs miss
structural equivalences (e.g., a refactored loop). Graph edit distance for
program dependence graphs is NP-hard. ZSS on ASTs provides a tractable,
semantically grounded distance \cite{Ding2024Revisiting}.

\subsection{KL-Matched Checkpoint Comparison}

\textbf{Rationale:} GRPO and DPO diverge from the base model at different rates per
training step. Step-aligned comparison confounds training speed with structural quality.
We instead compare checkpoints at matched KL divergence levels.

\textbf{KL matching procedure:}
\begin{enumerate}
  \item During training, log KL divergence at each checkpoint step $t$: $\kappa_t = D_\text{KL}(\pi_t \| \pi_\text{ref})$
  \item For each GRPO checkpoint $\pi_t^\text{GRPO}$ with $\kappa_t^\text{GRPO}$, find a DPO checkpoint $\pi_{t'}^\text{DPO}$ such that $|\kappa_{t'}^\text{DPO} - \kappa_t^\text{GRPO}| \leq \varepsilon$
  \item Form matched pairs $(\pi_t^\text{GRPO}, \pi_{t'}^\text{DPO})$ for statistical comparison
\end{enumerate}

We use tolerance $\varepsilon = 0.05$ ($\pm 5\%$), though
our preliminary run used $\varepsilon = 0.15$ due to limited checkpoint availability.

\textbf{Checkpoint diversity requirement:} KL matching requires sufficient checkpoint
density. We impose a pre-flight check: $|\{t : \text{checkpoint}_t \text{ is unique}\}| \geq N_\text{min}$
(we recommend $N_\text{min} = 10$) before proceeding to analysis. This check
prevents the checkpoint aliasing confound described in Section~\ref{sec:aliasing}.

\subsection{Statistical Testing}

For a collection of $K$ matched pairs $\{(\pi_k^\text{GRPO}, \pi_k^\text{DPO})\}_{k=1}^K$,
we compare SEP distributions using:

\textbf{Mann-Whitney U test} (primary): Tests whether $P(\text{SEP}^\text{GRPO} > \text{SEP}^\text{DPO}) > 0.5$
without assuming normality. Requires $K \geq 10$ unique pairs for adequate power.

\textbf{Bootstrap confidence interval} (secondary): 10,000 bootstrap resamples of
the mean SEP differential $\Delta = \text{mean}(\text{SEP}^\text{GRPO}) - \text{mean}(\text{SEP}^\text{DPO})$,
yielding a 95\% CI that does not rely on the matched-pairs assumption.

\textbf{Significance threshold:} $\alpha = 0.05$ for the Mann-Whitney test.

\subsection{Implementation}

The framework is implemented as a Python pipeline with five modules:

\begin{table}[h]
\centering
\caption{Framework module summary}
\label{tab:modules}
\begin{tabular}{ll}
\toprule
\textbf{Module} & \textbf{Function} \\
\midrule
\texttt{ast\_decomposition.py} & FA-AST node classification; SEP computation \\
\texttt{ast\_metric.py} & ZSS semantic edit distance (CF+DF nodes only) \\
\texttt{kl\_metric.py} & KL log computation; checkpoint matching \\
\texttt{sep\_analysis.py} & SEP analysis across matched checkpoint pairs \\
\texttt{statistical\_tests.py} & Mann-Whitney U, Spearman correlation, bootstrap CI \\
\bottomrule
\end{tabular}
\end{table}

\textbf{Training infrastructure:} We use TRL v1.3.0 \cite{vonWerra2020TRL}
GRPOTrainer and DPOTrainer on DeepSeek-Coder-7B-instruct-v1.5
\cite{Guo2024DeepSeek} with the following configuration:

\begin{table}[h]
\centering
\caption{Training hyperparameters}
\label{tab:hyperparams}
\begin{tabular}{lll}
\toprule
\textbf{Parameter} & \textbf{GRPO} & \textbf{DPO} \\
\midrule
Learning rate & 1e-6 & 5e-7 \\
Batch size & 4 & 2 \\
Gradient accumulation & 4 & 8 \\
KL penalty $\beta$ & 0.04 & 0.1 \\
Training steps & 1000 & 1000 \\
Checkpoint save interval & Every 100 steps & Every 100 steps \\
\bottomrule
\end{tabular}
\end{table}

Execution rewards use the \texttt{evalplus} harness \cite{Liu2023EvalPlus}:
binary (+1/0) and error-type rewards (categorized by exception class:
SyntaxError, RuntimeError, AssertionError, etc.).

\textbf{Evaluation dataset:} HumanEval+ (164 problems) and MBPP+ (378 problems)
from evalplus \cite{Liu2023EvalPlus}, covering standard function-level Python
code generation.
